% TOP_PROBLEM: {"id": 30006752, "title": "Uniform Subconvexity for Standard Automorphic L-Functions on GL(n)", "category_id": 1, "category": {"id": 1, "name": "number_theory", "display_name": "Number Theory", "description": "Properties of integers, prime numbers, Diophantine equations.", "slug": "number-theory", "order_index": 1, "created_at": "2026-07-31T15:26:25.670Z"}, "status": "open"}
% FIRST_POSTED: null
% !TeX program = lualatex
% Compile twice: lualatex 30006752.tex
% All references and prose are embedded; no BibTeX database or external inputs.
\documentclass[11pt,a4paper]{article}
\usepackage[margin=24mm,headheight=14pt]{geometry}
\usepackage{fontspec}
\setmainfont{Latin Modern Roman}
\setsansfont{Latin Modern Sans}
\setmonofont{Latin Modern Mono}
\usepackage{amsmath,amssymb,mathtools}
\usepackage{unicode-math}
\setmathfont{Latin Modern Math}
\usepackage{microtype}
\usepackage{enumitem,xurl,needspace}
\usepackage[dvipsnames]{xcolor}
\usepackage{hyperref}
\hypersetup{unicode=true,colorlinks=true,linkcolor=MidnightBlue,urlcolor=MidnightBlue,
 pdftitle={Research attempt: Problem 30006752},pdfauthor={Research notebook},bookmarksnumbered=true}
\usepackage{bookmark}
\usepackage{tocloft}
\setlength{\cftsecnumwidth}{3em}
\cftsetpnumwidth{2.5em}
\cftsetrmarg{4em}
\usepackage{titlesec}
\titleformat{\section}{\Large\bfseries\raggedright}{\thesection}{0.7em}{}
\usepackage{fancyhdr}
\pagestyle{fancy}\fancyhf{}
\fancyhead[L]{\small\sffamily Open Problems Math | Research notebook}
\fancyhead[R]{\small Problem 30006752 | 27 September 2026}
\fancyfoot[C]{\thepage}
\setlength{\parindent}{0pt}
\setlength{\parskip}{3pt plus 1pt}
\setlength{\emergencystretch}{3em}
\setcounter{tocdepth}{1}
\setcounter{secnumdepth}{1}
\setlist[itemize]{leftmargin=3em,itemsep=5pt,topsep=4pt,parsep=0pt}
\allowdisplaybreaks
\newcommand{\N}{\mathbb{N}}
\newcommand{\Z}{\mathbb{Z}}
\newcommand{\Q}{\mathbb{Q}}
\newcommand{\R}{\mathbb{R}}
\newcommand{\C}{\mathbb{C}}
\newcommand{\F}{\mathbb{F}}
\newcommand{\A}{\mathbb{A}}
\newcommand{\PP}{\mathbb P}
\newcommand{\E}{\mathbb{E}}
\newcommand{\eps}{\varepsilon}
\newcommand{\dd}{\,\mathrm d}
\newcommand{\sref}[2]{\hyperlink{source:#1:#2}{[S#2]}}
\newcommand{\eref}[2]{\hyperlink{extra:#1:#2}{[E#2]}}
% Turn the next switch off for a shorter reading copy. The quotations preserve
% every exactTarget field, including qualifications omitted from a short summary.
\newif\ifshowcatalogue
\showcataloguetrue
\newcommand{\cataloguescope}[1]{\ifshowcatalogue
 \paragraph{Exact scope in the supplied catalogue.}
 \begin{quote}\small #1\end{quote}\fi}
\usepackage{etoolbox}
\AtBeginEnvironment{itemize}{\small}
\numberwithin{equation}{section}
% Standalone export. Original research content preserved.
\begin{document}
\setcounter{section}{0}
% ================================================================
% problemId: 30006752
% Canonical problem ID: 30006752
% exactTarget SHA-256: ed2e33d61b5b05885eec9c52ad66159939d6d1f0abc4285439293d2cfb759a34
\clearpage
\section*{\texorpdfstring{Uniform Subconvexity for Standard Automorphic L-Functions on GL(n)}{Uniform Subconvexity for Standard Automorphic L-Functions on GL(n)}}\label{30006752}
\subsection{Definitions and mathematical statement}
For a unitary cuspidal $\pi$ on $\operatorname{GL}_n(\A_F)$, its analytic conductor $C(\pi)$ combines the finite conductor and the sizes of its archimedean gamma-factor parameters, in the cited normalization. The target is
\[
 \forall n,F\ \exists\delta_{n,F}>0,\ A_{n,F}>0\ \forall\pi:\qquad
 |L(\tfrac12,\pi)|\le A_{n,F}C(\pi)^{1/4-\delta_{n,F}}.
\]
The $L$-function is the finite part. All finite and archimedean parameters vary simultaneously, including unitary imaginary twists; neither a fixed representation family nor a single conductor aspect suffices. Any positive saving is allowed.
\par\smallskip\textit{Formulation and concept references:} \sref{30006752}{1}.
\cataloguescope{For every integer n\ensuremath{\ge}1 and every number field F, there exist constants \ensuremath{\delta}(n,F)>0 and A(n,F)>0 such that |L(1/2,\ensuremath{\pi})|\ensuremath{\le}A(n,F) C(\ensuremath{\pi})\textasciicircum{}(1/4\ensuremath{-}\ensuremath{\delta}(n,F)) for every unitary cuspidal automorphic representation \ensuremath{\pi} of GL\_n(A\_F), where L(s,\ensuremath{\pi}) is the standard finite-part L-function and C(\ensuremath{\pi}) its analytic conductor. The bound must be uniform simultaneously in all finite and archimedean parameters of \ensuremath{\pi}; imaginary unitary twists are included. Any fixed positive exponent saving suffices; the optimal exponent is not requested.}
\subsection{Short English statement}
Can one beat the convexity exponent uniformly for every standard automorphic $L$-function of each fixed degree over each number field?
% BEGIN RESEARCH ADDITION 147
\subsection{Research attempt}

\paragraph{Current frontier.}
The formulation source proves a $t$-aspect result for fixed representations
over $\Q$, and a spectral theorem under uniform growth of archimedean
parameters, with a polynomial finite-conductor loss. Its Section 1.2
explicitly distinguishes these from the universal-family target and records
the logarithmic general improvement and the full degree-two theorem.
\sref{30006752}{1} Hu--Nelson's 2025 Theorem 1.2 allows varying ramified sets but
retains specified local uniform-depth hypotheses and an uncontrolled-conductor
factor. \eref{30006752}{1} Berta's February 2026 manuscript treats prime-conductor
twists of a fixed degree-three representation over a number field, not every
parameter at once. \eref{30006752}{2}

\paragraph{Attack route.}
Try either to combine complementary aspect estimates or to control short
additive correlations of the actual approximate-functional-equation
coefficients. We calculate exactly when complementary estimates cover every
conductor ratio, then derive a quantitative sufficient correlation estimate.
The actual varying weights and full conductor remain visible throughout.

\paragraph{Attempt: exact aspect coverage.}
Let $C=QP$, with finite and archimedean conductor factors $Q,P\ge1$. Suppose,
hypothetically on the same family with uniform constants, that
\[
 |L|\ll Q^{1/4+a}P^{1/4-d_1},\qquad
 |L|\ll Q^{1/4-d_2}P^{1/4+b},
 \quad a,b\ge0,
\]
where $d_1,d_2>0$. With $x=\log Q/\log C$, the exponents relative to
$C^{1/4}$ are $ax-d_1(1-x)$ and $-d_2x+b(1-x)$. The maximum of their
minimum occurs at their intersection and equals
\begin{equation}\label{30006752:coverage}
 \frac{ab-d_1d_2}{a+b+d_1+d_2}.
\end{equation}
Consequently these bounds give a positive uniform saving exactly if
$ab<d_1d_2$. Otherwise the interval
\[
 \frac{d_1}{a+d_1}\le x\le\frac{b}{b+d_2}
\]
is uncovered, including a zero-saving boundary point when equality holds.
This describes what the bounds imply, not the true size of $L$ there.
For example $a=b=1/4$, $d_1=d_2=1/100$ leaves a large middle interval.
Taking a minimum or geometric mean does not repair it. The finite factor
in the source's Theorem 2 is not a harmless constant when finite conductor
varies. \sref{30006752}{1} A new bound must reduce those losses or cover the
middle regime, besides removing local parameter restrictions.

\paragraph{A finite weighted correlation inequality.}
Let $b_j$ be supported on $N$ consecutive integers and zero elsewhere. Set
\[
 S=\sum_jb_j,\quad Q_0=\sum_j|b_j|^2,\quad
 R_h=\sum_jb_{j+h}\overline{b_j}.
\]
For $1\le H\le N$,
\begin{equation}\label{30006752:vdc}
 |S|^2\le\frac{N+H-1}{H^2}
 \left(HQ_0+2\sum_{h=1}^{H-1}(H-h)\Re R_h\right).
\end{equation}
Indeed $HS=\sum_m\sum_{h=1}^H b_{m+h}$ and at most $N+H-1$ outer summands
are nonzero. Cauchy--Schwarz and expansion of the finite square give the
formula. Its bracket is nonnegative; no diagonal term has been discarded.

For a fixed $0<\eta\le1/2$, put $H=\lceil N^\eta\rceil$. Suppose
\begin{equation}\label{30006752:corr}
 Q_0\ll_\epsilon NC^\epsilon,\qquad
 \sum_{h=1}^{H-1}(H-h)\Re R_h
       \ll_\epsilon H^2N^{1-\eta}C^\epsilon.
\end{equation}
For $N\ge2$, $H\le N$, and \eqref{30006752:vdc} yields
\[
 |S|\ll_\epsilon C^{\epsilon/2}N^{1-\eta/2}.
\]
The case $N=1$ follows from the first estimate. Only an upper bound on the
weighted real sum is needed; bounding every absolute correlation is
stronger and could waste cancellation among shifts.

\paragraph{Conditional application retaining the true weights.}
Fix $n,F$. The uniform approximate functional equation supplies two
original/dual sums truncated at $C^{1/2+\tau}$ with rapidly decaying weights
depending on the actual gamma parameters, plus a rapidly decreasing error.
This is the uniform truncation input in Harcos's Theorem 1 and Corollary 1,
with its erratum, also recorded in his later dissertation, Section 1.2.
\eref{30006752}{3} No replacement by a parameter-independent cutoff is made.

Let $C=\max\{2,C(\pi)\}$ and aggregate Dirichlet coefficients $a_\pi(m)$
by ideal norm when $F\ne\Q$. On a dyadic interval $N\le m<2N$, set
\[
 b_{\pi,N}(m)=\mathbf1_{N\le m<2N}\mathbf1_{m\le C^{1/2+\tau}}
       (N/m)^{1/2}a_\pi(m)V_\pi(m/\sqrt C),
\]
using the genuine weight; do likewise for the dual sum. The precise missing
hypothesis is \eqref{30006752:corr} for these sequences, every $\epsilon>0$,
all allowed $\pi$ and all $N$, with one $\eta=\eta(n,F)>0$ and uniform
constants. The coefficient-square estimate is explicitly part of the
hypothesis, not an unproved Ramanujan assertion silently used below.

Each dyadic contribution is $N^{-1/2}S$, bounded by
$C^{\epsilon/2}N^{(1-\eta)/2}$. Since that exponent of $N$ is positive,
the dyadic sum is bounded geometrically by its largest scale. Choose
$\tau=\eta/16$, $\epsilon=\eta/32$. The final exponent is
\[
 \frac\epsilon2+\left(\frac12+\tau\right)\frac{1-\eta}{2}
 =\frac14-\frac\eta4+\frac{\eta(1-\eta)}{32}+\frac\eta{64}
 \le\frac14-\frac{13\eta}{64}<\frac14-\frac\eta8.
\]
Thus the correlation hypothesis implies the exact requested type of bound,
for example with $\delta_{n,F}=\eta/8$.

For pole characters in degree one, the contour shift retains its residue
terms rather than pretending the $L$-function is entire. Such characters
are norm twists of $\zeta_F$; with a Gaussian auxiliary factor their
residue contributions are
$O_F((1+|t|)^B e^{-t^2+O_F(|t|)})=O_F(1)$, and do not affect the deduction.
Alternatively these terms may simply be retained in the analytic input.
The two-sum formula without residue terms is used only in the entire case.

\paragraph{Stress test.}
For $b_j=1$ on $[1,N]$, $Q_0=N$ but $R_h=N-h$. With $H=o(N)$ the weighted
sum is asymptotic to $NH^2/2$, exceeding the proposed bound by a power.
Thus the diagonal estimate and functional-equation length do not generate
cancellation. Random mean-zero unit-modulus coefficients have expected
correlation zero, but are not automorphic evidence. A unitary twist
introduces phases $m^{-it}$ and does not guarantee uniform cancellation
in every short interval.

Keeping the true weight inside $b$ avoids assuming uniform partial summation
near degenerate archimedean parameters. Norm aggregation over $F$ is also
retained. The script checks the finite correlation identity, exact inequalities
for Gaussian-integer sequences and the aspect-coverage algebra, not the
unproved automorphic hypotheses.

\paragraph{Outcome.}
\textbf{Outcome: Conditional reduction.}
The attempt derives a two-aspect coverage threshold and a short-correlation
sufficient condition with explicit exponent bookkeeping. Proving that
condition uniformly in all ramification and archimedean regimes would yield
the target. Its arithmetic inputs remain unproved; no new universal
subconvexity theorem or illicit extension of a fixed-family result is claimed.

% END RESEARCH ADDITION 147
\subsection{Sources}
\begin{itemize}\raggedright
\item[\textnormal{[S1]}]\hypertarget{source:30006752:1}{} Bounds for standard L-functions.
\url{https://ultronozm.github.io/math/standard.html}.
{\small\textit{Catalogue formulation source.}}
% BEGIN RESEARCH REFERENCES 147
\item[\textnormal{[E1]}]\hypertarget{extra:30006752:1}{} Y. Hu and P. D. Nelson,
\emph{The subconvexity bound for standard L-function in level aspect},
arXiv:2503.12310v1, Sections 1.2--1.4, Theorem 1.2 and Corollary 1.3.
\url{https://arxiv.org/html/2503.12310v1}.
\item[\textnormal{[E2]}]\hypertarget{extra:30006752:2}{} F. Berta,
\emph{Subconvexity Problem on $\operatorname{GL}_3$ over number fields:
the twist aspect}, arXiv:2602.20095v1, 23 February 2026, stated main result.
\url{https://arxiv.org/abs/2602.20095}.
\item[\textnormal{[E3]}]\hypertarget{extra:30006752:3}{} G. Harcos,
\emph{Uniform approximate functional equation for principal L-functions},
IMRN 2002, 923--932, Theorem 1 and Corollary 1; erratum, IMRN 2004, 659--660.
\url{https://arxiv.org/abs/math/0111312}.
See also the author's \emph{Subconvex Bounds for Automorphic L-functions and
Applications}, 2011 dissertation, Section 1.2, whose bibliography includes
the erratum. \url{https://www.renyi.hu/~gharcos/ertekezes.pdf}.

% END RESEARCH REFERENCES 147
\end{itemize}

\end{document}
